\documentclass[11pt]{article}
\usepackage{amsmath, amssymb, amsthm}
\usepackage{algorithm}
\usepackage{algpseudocode}
\usepackage{geometry}
\usepackage{booktabs}
\usepackage{hyperref}
\usepackage{natbib}
\usepackage{graphicx}
\usepackage{xcolor}
\usepackage{todonotes}  % \todo{} and \missingfigure{} for tracking gaps

\geometry{margin=1in}

% --- Custom commands ---
\newcommand{\MG}{\mathcal{M}_G}
\newcommand{\EG}{\mathcal{E}_G}
\newcommand{\Itheta}{\mathcal{I}(\theta)}
\newcommand{\Lcf}{\mathcal{L}_{\text{cf}}}
\newcommand{\Ldist}{\mathcal{L}_{\text{dist}}}
\newcommand{\natgrad}{\widetilde{\nabla}}
\newcommand{\phil}{\phi_{ij}}
\newcommand{\thetaij}{\theta_{ij}}

% --- Theorem environments ---
\newtheorem{theorem}{Theorem}
\newtheorem{lemma}[theorem]{Lemma}
\newtheorem{proposition}[theorem]{Proposition}
\newtheorem{remark}{Remark}
\newtheorem{definition}{Definition}

% ================================================================
\title{Geodesic Attribution on the Graph Explanation Manifold:\\
  Navigating D4Explainer's Counterfactual Explanation Space\\
  via Fisher-Rao Geodesic Flow}
\author{Dylan L Holt\\
  \small Department of Applied and Computational Mathematics\\
  \small Johns Hopkins University\\
  \small \texttt{dholt14@jh.edu}}
\date{}
% ================================================================

\begin{document}

\maketitle

\begin{abstract}
% TODO: Fill in after results are complete.
% Draft seed:
We extend D4Explainer's discrete denoising diffusion framework for GNN
explainability by recasting its reverse diffusion process as geodesic
flow on the statistical manifold of graph distributions under the
Fisher-Rao metric. Where D4Explainer uses a Euclidean training objective
that ignores the curvature of the explanation manifold $\MG$, we propose
replacing it with natural gradient descent, enforcing that each
denoising step follows the geometry of the graph distribution space,
and yielding a canonically ordered, per-edge attribution decomposition
via martingale increments along the geodesic. We evaluate on the
D4Explainer benchmark suite (BA-Shapes, BA-Community, Tree-Cycles,
Tree-Grid) and report [results placeholder].
\todo[inline]{Add Results once ready.}
\end{abstract}
\setcounter{section}{0}
% ================================================================
\section{Introduction}\label{sec:paper}
% ================================================================

% --- Motivation ---
Graph Neural Networks (GNNs) have achieved strong performance across
node classification, graph classification, and link prediction tasks,
yet their lack of intrinsic interpretability limits deployment in
high-stakes settings \todo[inline]{Build upon GNN survey}.
Counterfactual explanations offer a principled alternative: given a
prediction $f(G)$, find a minimally edited graph $G'$ such that
$f(G') \neq f(G)$.

% --- Prior work gap ---
Early counterfactual methods such as CF-GNNExplainer
\todo[inline]{Recall CF-GNNExplainer} optimize edge masks in Euclidean space,
producing explanations that may lie off the data manifold. In other words, these explanations were not in the feasible range of the sample space, which was
a critical failure mode identified by \citet{doshi2017towards}.
D4Explainer \citep{chen2023d4explainer} addresses this by enforcing
the in-distribution constraint via discrete denoising diffusion,
representing a significant advance.

% --- The gap this paper fills ---
However, D4Explainer's denoising objective remains Euclidean: that is, the
optimizer treats the parameter space of edge probabilities as flat,
ignoring the curvature of the statistical manifold it implicitly
constructs. This motivates the central question of this work:
\emph{what is the geometrically correct way to navigate D4Explainer's
explanation space?}

% --- Contributions ---
We make three contributions:
\begin{enumerate}
  \item \textbf{Geometric reframing.} We show that D4Explainer's
    explanation space is a statistical manifold $\MG$ parameterized by
    Bernoulli edge probabilities, and that the natural metric on $\MG$
    is the Fisher-Rao metric $g_{ij}(\theta) =
    \delta_{ij}/(\thetaij(1-\thetaij))$.

  \item \textbf{Algorithmic replacement.} We replace D4Explainer's
    Euclidean gradient update with natural gradient descent, enforcing
    that each denoising step follows the geodesic on $(\MG,
    g_{\text{FR}})$.

  \item \textbf{Interpretive decomposition.} We derive per-edge
    attribution scores $\phil$ as martingale increments along the
    geodesic path, satisfying a completeness property analogous to
    Shapley efficiency: $\sum_{(i,j)} \phil =
    f(e^*_{\text{geo}}) - f(G)$.
\end{enumerate}

% ================================================================
\section{Background}
% ================================================================

\subsection{GNN Explainability}

\todo[inline]{Fill in Key Definitions for navigating paper: taxonomy of gradient-based, perturbation-based,
and counterfactual methods. Key papers: GNNExplainer, CF-GNNExplainer,
GraphSVX, Doshi-Velez \& Kim (2017), particular emphasis on the in-distribution
requirement as the motivating criterion.}

\subsection{Diffusion Models on Graphs}

Let $G = (V, E)$ be a graph with adjacency matrix
$A^{(0)} \in \{0,1\}^{n \times n}$.
The forward diffusion process corrupts $A^{(0)}$ over $T$ steps by
injecting noise, producing a sequence
$A^{(0)}, A^{(1)}, \ldots, A^{(T)}$
governed by the It\^{o} stochastic differential equation
\citep{oksendal2010}:
\begin{equation}
  dX_t = b(X_t)\,dt + \sigma(X_t)\,dW_t.
\end{equation}
The score-based reverse-time SDE \citep{song2021score} recovers the
data distribution by running:
\begin{equation}
  dX_t = \bigl[b(X_t) - \sigma^2 \nabla_x \log p_t(X_t)\bigr]\,dt
         + \sigma\,d\bar{W}_t,
\end{equation}
where $\nabla_x \log p_t$ is the score function --- the Fisher score
of $p_t$. This connection between reverse diffusion and the Fisher
score is the bridge to information geometry utilized in the proposed extension of D4Explainer.

\todo[inline]{Add: GDSS (Jo et al.\ 2022) as continuous-time graph
diffusion parent of D4Explainer. Connect to D4Explainer's discrete
adaptation and its loss $\mathcal{L} = \Ldist + \alpha \Lcf$ shown below.}
$$\\$$
The D4Explainer (Chen et al., 2023) utilizes the composite loss as follows:
\[
  \mathcal{L}(e_k) \;=\; \mathcal{L}_{\mathrm{dist}}(e_k) + \mathcal{L}_{\mathrm{cf}}(e_k),
\]
where $\mathcal{L}_{\mathrm{dist}}$ penalises deviation from $\mathcal{p}_{data}$ (measured by
Maximum Mean Discrepancy) and $\mathcal{L}_{\mathrm{cf}}$ penalises failure to flip the prediction. Maximum Mean Discrepancy (MMD) measures how much $e_k$ looks like a real graph from the training distribution (i.e., is $e_k$ a plausible in-distribution graph). D4Explainer trains $f_\theta$ with the Euclidean gradient $\nabla_\theta \mathcal{L}$,
treating the parameter space as flat and ignoring the curvature of $\MG$.

\subsection{Information Geometry}

\todo[inline]{Fill in Key Definitions: statistical manifold, Fisher-Rao metric,
natural gradient. Key papers: Amari (2016). State the key result:
natural gradient = Riemannian gradient on $(\mathcal{M}, g_\text{FR})$.
Also cite Ly et al.\ (2017) for Fisher information as operational metric.}

% ================================================================
\section{The Graph Explanation Manifold $\MG$}
% ================================================================

\begin{definition}[Graph Explanation Manifold]
Let $G = (V,E)$ be a graph. The \emph{graph explanation manifold}
$\MG$ is the statistical manifold parameterized by
$\theta \in (0,1)^{|E|}$, where each $\thetaij$ denotes the
probability that edge $(i,j)$ is present in a candidate explanation
$G'$. Each $\theta$ defines a distribution over graphs via independent
Bernoullis:
\begin{equation}
  p_\theta(A) = \prod_{(i,j) \in E}
    \thetaij^{A_{ij}}(1 - \thetaij)^{1 - A_{ij}}.
\end{equation}
\end{definition}

\begin{lemma}[Diagonal Fisher-Rao Metric on $\MG$]
Under the Bernoulli product parameterization, the Fisher-Rao metric
on $\MG$ is diagonal:
\begin{equation}
  g_{ij}(\theta) = \frac{1}{\thetaij(1 - \thetaij)}.
\end{equation}
\end{lemma}
\begin{proof}
  \todo[inline]{Move final proof to appendix - Direct computation from definition of Fisher information for
  Bernoulli family. From Ly et al.\ (2017), Section 2.}
\end{proof}

\begin{remark}
  The diagonal structure makes $\mathcal{I}(\theta)^{-1}$ computable
  in $O(|E|)$ rather than $O(|E|^2)$, which is the key computational
  advantage of the Bernoulli factorization.
\end{remark}

\paragraph{Independence assumption.}
The Bernoulli product parameterization treats edge probabilities as
independent. This is a modeling assumption shared with GDSS
\citep{jo2022score} and D4Explainer \citep{chen2023d4explainer}.
\todo[inline]{Cite GDSS and D4Explainer as
precedent since we are building under the same assumptions; flag relaxation to correlated edge model as future work.}

% ================================================================
\section{Geodesic Denoising via Natural Gradient}
% ================================================================

\subsection{D4Explainer Baseline (Algorithm~\ref{alg:d4})}

D4Explainer trains $f_\theta$ with the Euclidean gradient $\nabla_\theta \mathcal{L}$,
treating the parameter space as flat and ignoring the curvature of $\MG$.
This is the gap addressed in the proposed Algorithm 2, which implements a natural gradient on the explanation space versus the Euclidean gradient.

\begin{algorithm}[H]
\caption{D4Explainer Reverse Diffusion (Baseline)}
\label{alg:d4}
\begin{algorithmic}[1]
\Require $A^{(T)}$, denoiser $f_\theta$, classifier $f$, steps $T$, candidates $K$
\Ensure Best explanation $e^*$, candidate set $\{e_k\}$

\For{$k = 1$ \textbf{to} $K$}                     \Comment{Reverse diffusion}
  \State $A_k^{(T)} \leftarrow A^{(T)}$
  \For{$t = T$ \textbf{downto} $1$}
    \State $\widehat{A}_k^{(0)} \leftarrow f_\theta\!\left(A_k^{(t)}, t\right)$
    \Comment{Posterior mean: $\widehat{A}_k^{(0)} \approx
             \mathbb{E}\!\left[A^{(0)} \mid A^{(t)}\right]$}
    \State $A_k^{(t-1)} \sim p_\theta\!\left(A^{(t-1)} \mid
           A_k^{(t)},\, \widehat{A}_k^{(0)}\right)$
    \Comment{Sampling next step given current state and pred. clean graph}
    \State $A^{(t)}_t \leftarrow A_k^{(t-1)}$
    \Comment{Explicitly storing next step as input to next iteration}
  \EndFor
  \State $e_k \leftarrow \widehat{A}_k^{(0)}$
  \Comment{Store final predicted clean graph after $T$ steps as candidate explanation}
\EndFor

\For{$k = 1$ \textbf{to} $K$}                        \Comment{Evaluate loss}
  \State $\mathcal{L}_{\mathrm{dist}}(e_k)
         \leftarrow \mathrm{MMD}\!\left(e_k,\, \mathcal{p}_{data}\right)$
    \Comment{Measure deviation from in-distribution graphs}
  \State $\mathcal{L}_{\mathrm{cf}}(e_k)
         \leftarrow \mathrm{CrossEntropy}\!\left(f(e_k),\, 1 - \hat{Y}_G\right)$
    \Comment{Penalize if candidate does not flip prediction}
  \State $\mathcal{L}(e_k) \leftarrow
         \mathcal{L}_{\mathrm{dist}}(e_k) + \mathcal{L}_{\mathrm{cf}}(e_k)$
    \Comment{Sum for total loss of candidate explanation}
\EndFor

\State $e^* \leftarrow \operatorname*{arg\,min}_{e_k}\; \mathcal{L}(e_k)$
\Comment{Select candidate explanation with minimal total loss}
\State \Return $e^*,\; \{e_k\}_{k=1}^K$
\end{algorithmic}
\end{algorithm}

\todo[inline]{\textbf{Draw more from Key References:}
\newline\newline
Chen, J., Wu, S., Gupta, A., & Ying, R. (2023). D4Explainer: In-distribution GNN explanations via discrete denoising diffusion. Advances in Neural Information Processing Systems (NeurIPS 2023). arXiv:2310.19321.
\newline\newline
Jo, J., Lee, S., & Hwang, S. J. (2022). Score-based generative modeling of graphs via the system of stochastic differential equations. International Conference on Machine Learning, pp. 10362–10383. PMLR.}

\subsection{Proposed: Geodesic Denoising via Natural Gradient (Algorithm~\ref{alg:geo})}

\begin{lemma}[Natural Gradient as Riemannian Gradient]
The natural gradient update
$\theta \leftarrow \theta - \eta\,\Itheta^{-1}\nabla_\theta\mathcal{L}$
is the Riemannian gradient descent step on $(\MG, g_\text{FR})$.
\end{lemma}
\begin{proof}
  \todo[inline]{Move final proof to appendix - Recall Amari (2016), Chapter 3. The Riemannian gradient is
  $\text{grad}\,\mathcal{L} = G(\theta)^{-1}\nabla\mathcal{L}$ where
  $G(\theta)$ is the metric tensor. For Fisher-Rao, $G = \Itheta$.}
\end{proof}

D4Explainer's reverse diffusion optimizes a Euclidean gradient
$\nabla_\theta \mathcal{L}$, treating the parameter space as flat and
ignoring the curvature of the graph distribution manifold $\mathcal{M}_G$.
We replace this with a \emph{natural gradient} update that follows the
geodesic on $(\mathcal{M}_G, g_{\mathrm{FR}})$, enforcing that each
denoising step respects the local geometry of the edge-probability
distributions rather than moving uniformly in adjacency matrix space.
\newline\newline
For the product Bernoulli edge model (from Algorithm 1), the Fisher information matrix is
diagonal with entries
$$
    \mathcal{I}(\theta)_{(i,j),(i,j)}
    \;=\;
    \frac{1}{\theta_{ij}\,(1 - \theta_{ij})},
$$
which diverges as $\theta_{ij} \to \{0, 1\}$ (high curvature near
deterministic edges) and attains its minimum of $4$ at
$\theta_{ij} = 0.5$ (flat region, maximally uncertain edges).
Inverting this metric yields the natural gradient
$$
    \widetilde{\nabla}_\theta \mathcal{L}[ij]
    \;=\;
    \theta_{ij}(1 - \theta_{ij})\cdot \nabla_\theta \mathcal{L}[ij],
$$
where the scaling factor $\theta_{ij}(1 - \theta_{ij})$ is a concave
function of $\theta_{ij}$, maximized at $0.5$ and vanishing at the
boundaries. The geodesic update is then
$$
    \theta^{(t-1)}
    \;=\;
    \theta^{(t)}
    \;-\;
    a_t\,\mathcal{I}(\theta^{(t)})^{-1}\,\nabla_\theta \mathcal{L},
$$
where $\{a_t\}= \frac{1}{t}$ such that early steps are large but become smaller at each iteration.
The elementwise form makes the geometric content transparent: edges with
$\theta_{ij} \approx 0.5$ (uncertain) receive large updates and resolve
quickly, while edges near $0$ or $1$ (decided) are protected from large
perturbations. This means that noisy edges are resolved more quickly while preserving edges with higher certainty. The method follows these natural gradients from step to step, which is our proposed extension.

\begin{algorithm}[H]
\caption{Geodesic Denoising via Fisher-Rao Natural Gradient (Proposed)}
\begin{algorithmic}[1]
\Require $A^{(T)}$, classifier $f$, steps $T$,
         Fisher oracle $\mathcal{I}(\cdot)$, step sizes $\{a_t\}$, threshold $\tau$
\Ensure Geodesic path $\{\theta^{(t)}\}_{t=0}^T$, explanation $e^*_{\mathrm{geo}}$

\State $\theta^{(T)}[ij] \leftarrow A^{(T)}[ij]$
\Comment{Initialise at noisy graph; $\theta \in \{0,1\}^{|E|}$}

\For{$t = T$ \textbf{downto} $1$}

  \State $\nabla\mathcal{L} \leftarrow
         \nabla_\theta\!\left[\mathcal{L}_{\mathrm{dist}}\!\left(\theta^{(t)}\right)
         + \mathcal{L}_{\mathrm{cf}}\!\left(\theta^{(t)}, f\right)\right]$
  \Comment{Euclidean gradient of composite loss (same as D4Explainer)}

  \State $\mathcal{I}_t \leftarrow
         \mathcal{I}\!\left(\theta^{(t)}\right)$
  \Comment{Fisher information at current manifold point}

  \State $\natgrad\mathcal{L} \leftarrow \mathcal{I}_t^{-1}\,\nabla\mathcal{L}$
  \Comment{Natural gradient adj. from Euclidean gradient: steepest descent on $(\MG, g_{\mathrm{FR}})$}

  \State $\theta^{(t-1)} \leftarrow \theta^{(t)} - a_t\,\natgrad\mathcal{L}$
  \Comment{Geodesic step along $\MG$ with step-size $a_t=\frac{1}{t}$}

  \State $\theta^{(t-1)}[ij] \leftarrow
         \mathrm{Clip}\!\left(\theta^{(t-1)}[ij],\,\varepsilon,\,1-\varepsilon\right)$
  \Comment{Avoid boundary degeneracy where $\mathcal{I}$ diverges}

  \State Store $\left(\theta^{(t-1)},\, \mathcal{I}_t,\,
         \natgrad\mathcal{L}\right)$ in $\mathrm{geodesic\_path}$

\EndFor

\State $e^*_{\mathrm{geo}}[ij]
    \;=\;
    \mathbf{1}\!\left[\,\theta^{(0)}_{ij} > \tau\,\right]$
\Comment{Adjacency matrix computed from edge probabilities against threshold $\tau$}

\State \Return $e^*_{\mathrm{geo}},\; \mathrm{geodesic\_path}$
\end{algorithmic}
\end{algorithm}

\todo[inline]{\textbf{Draw more from Key References:}
\newline\newline
Amari, S. (2016). Information geometry and its applications. Springer.
\newline\newline
Coifman, R. R., & Lafon, S. (2006). Diffusion maps. Applied and Computational Harmonic Analysis, 21(1), 5–30.
\newline\newline
Song, Y., Sohl-Dickstein, J., Kingma, D. P., Kumar, A., Ermon, S., & Poole, B. (2021). Score-based generative modeling through stochastic differential equations. ICLR 2021. arXiv:2011.13456.}

\paragraph{Implementation note.}
The natural gradient is applied as a backward hook on
\texttt{score\_batch} before \texttt{optimizer.step()}, leaving
the loss functions \texttt{loss\_dist} and \texttt{loss\_cf} unchanged:
\begin{verbatim}
score_batch = model(A=..., node_features=..., mask=..., noiselevel=sigma)
register_natural_gradient_hook(score_batch)   # hook fires in backward()
...
loss = loss_dist + args.alpha_cf * loss_cf
loss.backward()      # natural gradient rescaling applied here
optimizer.step()     # steps in corrected direction
\end{verbatim}

\todo[inline]{Add Github Repository reference}

% ================================================================
\section{Attribution Decomposition via Martingale Increments}
% ================================================================

The geodesic path $\{\theta^{(t)}\}_{t=T}^{0}$ induces a martingale
along the GNN prediction:
\begin{equation}
  M_t = \mathbb{E}\bigl[f(e^*_{\text{geo}}) \mid \theta^{(t)}\bigr],
  \quad M_T = f(G), \quad M_0 = f(e^*_{\text{geo}}).
\end{equation}
Per-edge attribution scores are defined as:
\begin{equation}
  \phil = \sum_{t=0}^{T-1} \Delta M_t \cdot
    \frac{|\theta^{(t+1)}_{ij} - \theta^{(t)}_{ij}|}
         {\sum_{(a,b)}|\theta^{(t+1)}_{ab} - \theta^{(t)}_{ab}|}.
\end{equation}

\begin{algorithm}[H]
\caption{Attribution Decomposition via Geodesic Martingale Increments}
\label{alg:attr}
\begin{algorithmic}[1]
\Require Geodesic path $\{\theta^{(t)}\}$, classifier $f$,
         input $G$, explanation $e^*_{\text{geo}}$
\Ensure Attribution scores $\{\phil\}$, ranked edge list
\For{$t = 0$ \textbf{to} $T$}
  \State $M_t \leftarrow f(\theta^{(t)})$
\EndFor
\For{each edge $(i,j) \in E$}
  \State $\phil \leftarrow 0$
  \For{$t = 0$ \textbf{to} $T-1$}
    \State $\Delta M_t \leftarrow M_{t+1} - M_t$
    \State $w_{ij}^{(t)} \leftarrow
      \dfrac{|\theta^{(t+1)}_{ij} - \theta^{(t)}_{ij}|}
            {\sum_{(a,b)}|\theta^{(t+1)}_{ab} - \theta^{(t)}_{ab}|}$
    \State $\phil \leftarrow \phil + \Delta M_t \cdot w_{ij}^{(t)}$
  \EndFor
\EndFor
\State \textbf{Assert}\;
  $\textstyle\sum_{(i,j)}\phil \approx f(e^*_{\text{geo}}) - f(G)$
  \Comment{Completeness}
\State \Return $\{\phil\}$,\;
  $\text{ranked} \leftarrow \text{SortDesc}(\{\phil\})$
\end{algorithmic}
\end{algorithm}

\begin{proposition}[Completeness]
The attribution scores satisfy:
\begin{equation}
  \sum_{(i,j) \in E} \phil = f(e^*_{\text{geo}}) - f(G).
\end{equation}
\end{proposition}
\begin{proof}
  \todo[inline]{Move final proof to appendix - Show completeness follows from telescoping the martingale increments and the normalization of edge weights $w_{ij}^{(t)}$. Note: this holds exactly in continuous time; in the discrete approximation, state the error bound or frame as approximate.}
\end{proof}

\todo[inline]{Address Shapley analogy to cement comparison in familiar explainability framework: completeness is in
the spirit of Shapley efficiency but is not a formal equivalence.
Compare to GraphSVX (Duval \& Malliaros, 2021)}

% ================================================================
\section{Experiments}
% ================================================================

\subsection{Setup}

\paragraph{Datasets.}
Following \citet{chen2023d4explainer}, we evaluate on the four benchmark datasets for Node Classification (BA-Shapes, Tree-Cycles, Tree-Grid, and Cornell's webpage graph) and four benchmark datasets for Graph Classification (BA-3Motif, Mutag, BBBP, and NCI1). 
\todo[inline]{Pull from Chen et al. 2023 - Add dataset statistics table: nodes, edges, classes, train/test split, and GNN classifier for each problem; Put details on training parameters in appendix.}
\todo[inline]{If computation becomes a challenge, focus on Graph Classification as true scope}

\paragraph{Baselines.}
\begin{itemize}
  \item \textbf{D4Explainer} \citep{chen2023d4explainer}:
    primary baseline; same initialization, same GNN backbone (GCN).
  \item \textbf{GNNExplainer}: gradient-based mask (Add reference).
  \item \textbf{CF-GNNExplainer}: Euclidean counterfactual (Add reference).
\end{itemize}

\paragraph{Metrics (following \citet{chen2023d4explainer}).}
\begin{center}
\begin{tabular}{lll}
\toprule
Metric & Definition & Target \\
\midrule
CF-ACC    & Fraction of explanations flipping prediction & Higher \\
Proximity & Avg.\ edit distance $d(G, G')$ & Lower \\
Validity  & In-distribution score (discriminator) & Higher \\
Sparsity  & Fraction of edges retained in $G'$ & Lower \\
\bottomrule
\end{tabular}
\end{center}

\todo[inline]{Add formal computations for each metric here, following pattern from Chen et al. (2023) paper}

\paragraph{Geometric validation (additional).}
\todo[inline]{Add one experiment measuring discriminator/FID score
along the denoising path for D4Explainer vs.\ geodesic approach.
This is the key experiment that validates the geometric claim
directly, not just the downstream metrics.}

\subsection{Hyperparameter Study}

\todo[inline]{Report: $\tau$ (threshold) vs.\ CF-ACC/Sparsity
Pareto frontier; step size $\eta$; ablation over $T$ (diffusion
steps).}

\subsection{In-Distribution Evaluation}

\todo[inline]{Following Chen et al. (2023) paper, add MMD results for same (3) select real world datasets; Consider work from Theory of Stats II project; Report and compare to D4Explainer as additional metric}

% \subsection{Ablation}

% \todo[inline]{Two key ablations:
% (1) Natural gradient vs.\ standard gradient (isolate geometric
%     contribution).
% (2) With vs.\ without attribution decomposition (verify completeness
%     property empirically).}

\subsection{Results}

\todo[inline]{Insert results table here following D4Explainer format.
Placeholder:}

\begin{center}
\begin{tabular}{lcccc}
\toprule
Method & CF-ACC $\uparrow$ & Proximity $\downarrow$
       & Validity $\uparrow$ & Sparsity $\downarrow$ \\
\midrule
GNNExplainer      & --- & --- & --- & --- \\
CF-GNNExplainer   & --- & --- & --- & --- \\
D4Explainer       & --- & --- & --- & --- \\
\textbf{Ours}     & --- & --- & --- & --- \\
\bottomrule
\end{tabular}
\end{center}

% \missingfigure{Geodesic path visualization on BA-Shapes: edge
% probability trajectories $\theta^{(t)}$ vs.\ D4Explainer path,
% overlaid on manifold schematic.}

% ================================================================
\section{Discussion}
% ================================================================

\subsection{Connection to Prior Work}

\todo[inline]{Compare to: GraphSVX (Shapley-based, sampling
approximation); GNNExplainer (gradient masking, no distributional
constraint); D4Explainer (our direct baseline). Note that our
completeness property holds by construction vs.\ approximation in
GraphSVX.}

\todo[inline]{Note connection to Coifman \& Lafon (2006): diffusion
maps provide an independent spectral characterization of $\MG$ via
the eigenfunctions of the graph Laplacian. This is a natural future
extension (see Section~\ref{sec:future}).}

\subsection{Stochastic Approximation Interpretation}

D4Explainer's training loop instantiates stochastic approximation
in the Robbins-Monro sense. The denoising model is trained by
minimizing $\mathcal{L} = \Ldist + \alpha\Lcf$; because the forward
diffusion injects randomness at each timestep, every gradient
estimate is computed on a stochastically corrupted input, making the
training update a genuine SA recursion:
\begin{equation}
  \theta_{n+1} = \theta_n - a_n\,\hat{g}(\theta_n, \xi_n),
\end{equation}
where $\xi_n$ encodes both mini-batch randomness and diffusion noise.
\todo[inline]{Recall Robbins \& Monro (1951). Note: this framing is not used
in the D4Explainer paper itself but is a perspective for the proposed approach.}

\subsection{Limitations}

\begin{itemize}
  \item \textbf{Bernoulli independence.} Edges in real graphs are
    not independent; the diagonal FIM is an approximation.
  \item \textbf{Martingale discretization.} The martingale structure
    holds exactly in continuous time; the discrete approximation
    introduces error that should be bounded in future work.
  \item \textbf{Scalability.} The diagonal FIM is $O(|E|)$; extension
    to correlated edge models would require $O(|E|^2)$ computation
    or structured approximation.
\end{itemize}

% ================================================================
\section{Conclusion}
% ================================================================

\todo[inline]{Write after results. Restate three contributions;
position as opening of a research program in geometrically principled
GNN interpretability.}

% ================================================================
\section{Future Work and Open Questions}
\label{sec:future}
% ================================================================
\todo[inline]{Tracking several notes for this section based on Optimal Stopping and other course modules.}
% The following questions were identified during this work but not
% addressed within its scope. They are recorded here as a research
% agenda.

% \paragraph{Label-preserving forward diffusion.}
% D4Explainer does not guarantee that $f(A^{(T)}) = f(A^{(0)})$
% during forward diffusion. The Fisher-Rao geometry on $\MG$ could
% provide a principled label-preserving forward schedule by constraining
% the diffusion path to remain within a geodesic ball around the
% original prediction. This is the strongest open question arising
% from this work.

% \paragraph{Correlated edge models.}
% Relaxing the Bernoulli independence assumption to a structured
% graphical model (e.g.\ a Markov random field over edges) would
% require a non-diagonal FIM and a more general natural gradient
% approximation (e.g.\ K-FAC style). This connects to the broader
% question of which graph generative models are compatible with
% tractable information geometry.

% \paragraph{Diffusion maps on $\MG$.}
% \citet{coifman2006diffusion} show that diffusion maps provide a
% spectral characterization of a manifold via the eigenfunctions of
% the graph Laplacian operator. Applying this to $\MG$ would yield
% intrinsic coordinates for the explanation manifold, providing an
% independent validation of the geodesic paths computed here.

% \paragraph{Shapley axiom formalization.}
% The completeness property $\sum_{(i,j)}\phil = f(e^*) - f(G)$
% is in the spirit of Shapley efficiency but is not a formal
% equivalence. A full axiomatic treatment --- showing which Shapley
% axioms the geodesic attribution satisfies and which it does not ---
% would sharpen the comparison to GraphSVX \citep{duval2021graphsvx}
% and GNNShap.

% \paragraph{Extension to node-level tasks.}
% The current framework addresses graph-level classification.
% Extending geodesic attribution to node-level tasks would require
% adapting the manifold parameterization to account for the
% neighborhood structure of the target node.

% ================================================================
% Bibliography
% ================================================================
\bibliographystyle{plainnat}
\begin{thebibliography}{99}

\bibitem[Amari(2016)]{amari2016}
Amari, S. (2016).
\textit{Information geometry and its applications}.
Springer.

\bibitem[Chen et al.(2023)]{chen2023d4explainer}
Chen, J., Wu, S., Gupta, A., \& Ying, R. (2023).
D4Explainer: In-distribution GNN explanations via discrete denoising
diffusion.
\textit{Advances in Neural Information Processing Systems (NeurIPS 2023)}.
arXiv:2310.19321.

\bibitem[Coifman \& Lafon(2006)]{coifman2006diffusion}
Coifman, R.~R., \& Lafon, S. (2006).
Diffusion maps.
\textit{Applied and Computational Harmonic Analysis}, 21(1), 5--30.

\bibitem[Doshi-Velez \& Kim(2017)]{doshi2017towards}
Doshi-Velez, F., \& Kim, B. (2017).
Towards a rigorous science of interpretable machine learning.
arXiv:1702.08608.

\bibitem[Duval \& Malliaros(2021)]{duval2021graphsvx}
Duval, A., \& Malliaros, F. (2021).
GraphSVX: Shapley value explanations for graph neural networks.
\textit{European Conference on Machine Learning (ECML 2021)}.

\bibitem[Jo et al.(2022)]{jo2022score}
Jo, J., Lee, S., \& Hwang, S.~J. (2022).
Score-based generative modeling of graphs via the system of stochastic
differential equations.
\textit{International Conference on Machine Learning (ICML 2022)},
pp.~10362--10383.

\bibitem[Ly et al.(2017)]{ly2017tutorial}
Ly, A., Marsman, M., Verhagen, J., Grasman, R.~P., \& Wagenmakers,
E.~J. (2017).
A tutorial on Fisher information.
\textit{Journal of Mathematical Psychology}, 80, 40--55.

\bibitem[{\O}ksendal(2010)]{oksendal2010}
{\O}ksendal, B. (2010).
\textit{Stochastic differential equations: An introduction with
applications} (6th ed.).
Springer.

\bibitem[Song et al.(2021)]{song2021score}
Song, Y., Sohl-Dickstein, J., Kingma, D.~P., Kumar, A., Ermon, S.,
\& Poole, B. (2021).
Score-based generative modeling through stochastic differential
equations.
\textit{International Conference on Learning Representations (ICLR 2021)}.

\end{thebibliography}

\end{document}